\documentclass[10pt,a4paper]{amsart}
\usepackage[margin=19mm]{geometry}
\usepackage{amsmath,amssymb,mathtools}
\usepackage[scr=rsfs,cal=euler]{mathalfa}
\usepackage{xfrac}
\usepackage[unicode,hidelinks,bookmarks=false]{hyperref}
\newcommand{\ie}{i.e.\ }
\newcommand{\cf}{cf.\ }
\newcommand{\resp}{respectively\ }
\newcommand{\Subsection}{\subsection}
\providecommand{\nobreakdash}{\nobreak\hbox{-}}
\setlength{\emergencystretch}{2em}
\multlinegap0pt
\usepackage{amssymb}
\usepackage{mathtools}
\usepackage{bm}
\usepackage{enumerate}
\usepackage{tikz}
\usepackage{xcolor}
\usepackage{float}
\usepackage[shortlabels]{enumitem}
\newtheorem{theorem}{Theorem}
\newtheorem{corollary}{Corollary}
\title{On roots of domination polynomials for friendship and book graphs}
\author{Bilal Ahmad Rather}
\date{Source: arXiv:2604.08998v2 (CC BY 4.0)}
\begin{document}
\maketitle

\noindent\emph{Source and licence.} On roots of domination polynomials for friendship and book graphs. Version arXiv:2604.08998v2 (CC BY 4.0), distributed by arXiv under the Creative Commons Attribution 4.0 International licence, http://creativecommons.org/licenses/by/4.0/, as recorded in the arXiv licence metadata for this exact version and shown on the abstract page https://arxiv.org/abs/2604.08998v2. Full licence notice: \texttt{licenses/arxiv-2604.08998-CC-BY-4.0.txt}.\par\smallskip

\noindent\textit{Editorial scope.} Counted unit: the article's corollary cor:monotone (source0.tex lines 291--300). The article's complete original proof of it is delivered with it (source0.tex lines 302--328, one \texttt{proof} environment of 27 lines). No mathematics is written, repaired or re-derived: the statement and the proof are the article's own lines, in the article's own notation, and no retained line is substituted or re-flowed.  Retained uncounted as the source-local context the proof needs: lines 224--229.  Nothing was omitted from any retained span, so the package carries no omission record.  No retained line contains a citation, so the wrapper carries no bibliography and no bibliography entry.  No retained line contains a figure environment, an image-inclusion command or drawing code, so no image is delivered and the package declares no asset dependency.  Editorial formatting only, in addition: an AMS \texttt{amsart} preamble that restates the article's own declarations and macros used by the retained lines, namely the environments \texttt{theorem}, \texttt{corollary} on separate counters, together with the standard packages those lines need.  The retained environments print in this document as Theorem 1, Corollary 1 in order of appearance, whereas the article prints the same items with the numbering of its own full document.  The complete native source of the article as distributed in the arXiv e-print for this exact version is bundled unchanged as \texttt{sources/198184/source0.tex}.
	\begin{theorem} \label{thm:exactthree}
		If $n$ is even, then $D(F_n,x)$ has exactly three real zeros, namely $0$ and two simple zeros
		$$
		x_n^- \in (-2,-1), \qquad x_n^+ \in (-1,0).
		$$
	\end{theorem}

	\begin{corollary}
		\label{cor:monotone}
		For even $n$, the roots from Theorem~\ref{thm:exactthree} satisfy
		$$
		x_n^- \searrow -1-\frac{1}{\sqrt2},
		\qquad
		x_n^+ \searrow -1+\frac{1}{\sqrt2}
		$$
		along the subsequence of even integers $n=2,4,6,\dots$.
	\end{corollary}

	\begin{proof}
		Let $t_n=x_n^++1\in(0,1)$ and $s_n=-(x_n^-+1)\in(0,1)$. Then from the proof of Theorem~\ref{thm:exactthree}, we have
		$$
		(1-t_n)^{n-1}(1+t_n)^n=t_n^{2n},\quad \text{and} \quad 	(1-s_n)^n(1+s_n)^{n-1}=s_n^{2n}.
		$$
		 Clearly, we see that $t_n>\tfrac{1}{\sqrt{2}}$. The initial identity shows that \(\left(\frac{t_n^2}{1-t_n^2}\right)^n=\frac{1}{1-t_n}>1 \), therefore \(\frac{t_n^2}{1-t_n^2}>1 \), implying that \(t_n>\frac{1}{\sqrt{2}} \).  From Theorem \ref{thm:exactthree}, the difference of functions $\phi_n$ and $\phi_{n+2}$ is
		 \begin{equation*}
		 	\phi_{n+2}(t)-\phi_n(t)=2\log\left(\frac{t^2}{1-t^2}\right).
		 \end{equation*}
		 At \(t=t_n \), the difference \(\phi_{n+2}(t)-\phi_n(t) \) is positive, indicating that \(t_n > \frac{1}{\sqrt{2}} \). Because $\phi_{n+2}$ is absolutely expanding and $\phi_{n+2}(t_n) > 0$, the zero must be to the left of $t_n$. Thus, $t_{n+2} < t_n$, implying $x_{n+2}^+ < x_n^+$.
		With a similar idea, the second identity implies that $s_n<\tfrac{1}{\sqrt2}$, since
		$$
		\left(\frac{s_n^2}{1-s_n^2}\right)^n=\frac{1}{1+s_n}<1.
		$$
		Theorem \ref{thm:exactthree} implies that
		$$
		\psi_{n+2}(s)-\psi_n(s)=2\log\left(\frac{s^2}{1-s^2}\right).
		$$
		At $s=s_n$, the value $\psi_{n+2}(s)-\psi_n(s)$ is negative,  so $\psi_{n+2}(s_n)<0$, and since $\psi_{n+2}$ is strictly rising, its zero is to the right of $s_n$. Thus, $s_{n+2}>s_n$, and it implies that $x_{n+2}^-<x_n^-$. The sequences ${t_n}$ and ${s_n}$ are monotone, bounded, and hence convergent. To evaluate the limit in $\left(\frac{t_n^2}{1-t_n^2}\right)^n=\frac{1}{1-t_n}$, the only solution is $\tfrac{t^2}{1-t^2}=1$, which gives $t=\tfrac{1}{\sqrt2}$.
		 Likewise, $s_n\to \tfrac{1}{\sqrt2}$. Thus, we get
		$$
		x_n^+=t_n-1 \to -1+\frac{1}{\sqrt2},
		\qquad\text{and} \qquad
		x_n^-=-1-s_n \to -1-\frac{1}{\sqrt2}.
		$$
		The demonstrated monotonicity indicates that the convergence is decreasing in both scenarios.
	\end{proof}

\end{document}
